
\section{Introduction} \label{s:intro}
Databases have proven to be a useful and versatile container for housing structured data, consisting of relatively simple but well-defined data types, referred to as \emph{scalars} (mostly numbers and strings).  In the 1970's databases evolved beyond a basic storage container to provide functions critical to the business by allowing relations amongst the stored data to be expressed and persisted \cite{Codd:1970:RMD:362384.362685}. This is fundamental to how business operates---where data from one part of the business must be correlated to other parts of the business if insights and efficiencies are to be realized.

As a simple example, consider operations at a retailer such as Walmart where a database has two tables: one records point-of-sale transactions, the other holds inventory data. At the point-of-sale it's useful to have a record of all the items purchased, however it's far more useful to be able to then relate this back to the inventory system to automatically reorder goods as needed. This relation, and the automation it enabled, let Walmart eliminate manual shelf stock planning, lowered cost, and prevented over- and under-stocks. This basic example shows where adding relation provides significant efficiency to the business.

The container used to store such transactions is typically a relational \db{}, comprised of rows (tuples) and columns of scalar data. A substantial strength of relational databases comes in the form of a standard query and transformation logic, SQL, which allows applications outside the creating \app{} to query the data \cite{Angles:2008:SGD:1322432.1322433}. This ability to separate data from the bounds of the creating \app{} amplifies and extends the value the data can provide the business.

\subsection{Database for Unstructured Data} \label{s:uddb}
In contrast to the simple scalar types typical to \sd{}, \ud{} is comprised of rich and expressive data types (e.g. PowerPoint presentations or full-motion video) that do not fit well into the traditional database paradigm.

An \os{} is essentially a database for unstructured data, composed of two parts: a distributed database that holds \emph{object references} and a distributed file store that stores the user data, referred to as \emph{data objects} or \emph{blobs}.  The database is typically modeled as a NoSQL ``shared nothing'' data store for horizontal scaling---replicating and partitioning data over many servers \cite{Cattell:2011:SSN:1978915.1978919}.

Public cloud examples include key-value stores such as Amazon's Dynamo \cite{DeCandia:2007:DAH:1323293.1294281} and Project Voldemort used by LinkedIn \cite{Sumbaly:2012:SLB:2208461.2208479}. In the enterprise and service provider sectors, the Hitachi Content Platform\footnote{\href{http://www.hds.com/products/file-and-content/content-platform}{http://www.hds.com/products/file-and-content/content-platform}}  (HCP) provides a distributed object store that typically resides behind a firewall \cite{RJP:2010:DOS-POP}. HCP is conceptually similar to the combination of Google's Bigtable \cite{Chang:2008:BDS:1365815.1365816} for storing object references and Google File System (GFS) \cite{Ghemawat:2003:GFS:945445.945450} for storing data objects.

\subsection{Abstraction}
Fundamental to object storage is that the detail of the distributed database and underlying file systems are abstracted from both the client \apps{} (users) and system administrators. In the process \oss{} shift the client model, essentially presenting storage as a service rather than requiring clients to be directly involved in data storage decisions, such as properly balancing directory trees. While some \oss{} only support a single flat namespace \cite{RJP:2005:CA}, others allow the global namespace to be logically partitioned for greater security, e.g. with a collection of \emph{buckets} in Amazon S3 \cite{varia2008cloud}, or \emph{namespaces} in HCP \cite{MR:2013:HCP-ARCH}.

Such abstraction allows for comparatively na\"{\i}ve users and administrators as the \os{} takes care of the detail of \emph{how} and \emph{where} user data is stored, protected, geo-replicated,  de-duplicated, versioned, garbage collected, and so forth. This grossly simplifies application development and deployment, as evidenced by the plethora of start-ups who are able in short order to get a worldwide service up, running, and generating revenue by leveraging a public object storage service---something never seen before in history.

Likewise, we see where limited compute platforms, such as smartphones and tablets, are able to have full access to a universe of data despite zero support for a single storage protocol. Such devices have no support for traditional storage protocols (FibreChannel, SCSI, NFS, CIFS) made popular in the last century for LAN-attached disk \cite{gibson1998cost}. Instead both simple and complex data types are served by nothing more than the basic web protocol HTTP.

\subsection{Contributions}
This paper describes a mechanism for adding a \emph{relational layer} on top of the object storage layer, without destroying the simplicity gained through the abstractions for which \oss{} are noted. Our goal is for users to be able to manage relations between data in similar fashion to a relational database. However, since unstructured data is intrinsically different from structured data, it's necessary that we provide a substrate for defining and describing such relation. Further, we envision creating a mechanism by which users can query an RDOS in a manner similar to the way they query a relational database today for similar benefit.

\subsection{Document Organization}
The remainder of the document is organized as follows. Section \ref{s:ud} defines \ud{} and touches on some of the associated challenges this datatype brings. Metadata is essential to the challenge of creating a relational layer on top of \ud{}. Section \ref{s:md} defines \md{}, explains why it is of such value, and identifies its potential to transform \ud{} from an undifferentiated data mass to a highly correlated data set that provides genuine value to the business. Section \ref{s:hcp} describes HCP in its present form and section \ref{s:rdos} describes the ideal of extending it to be a \RDOS{} and the business value it provides. Related work is described in section \ref{s:relwork}  and section \ref{s:conclusion} concludes with a summary of the topics covered in this paper.
